\documentclass[10pt,a4paper]{amsart}
\usepackage[margin=19mm]{geometry}
\usepackage{amsmath,amssymb,mathtools}
\usepackage[scr=rsfs,cal=euler]{mathalfa}
\usepackage{xfrac}
\usepackage[unicode,hidelinks,bookmarks=false]{hyperref}
\newcommand{\ie}{i.e.\ }
\newcommand{\cf}{cf.\ }
\newcommand{\resp}{respectively\ }
\newcommand{\Subsection}{\subsection}
\providecommand{\nobreakdash}{\nobreak\hbox{-}}
\setlength{\emergencystretch}{2em}
\multlinegap0pt
\usepackage{amssymb}
\usepackage{mathtools}
\usepackage[shortlabels]{enumitem}
\newtheorem{theorem}{Theorem}
\usepackage{cleveref}
\crefname{theorem}{Theorem}{Theorems}
\newcommand{\down}{\mathord{\downarrow}}
\newcommand{\leT}{\leq_T}
\newcommand{\lfp}{\operatorname{lfp}}
\title{\textbf{Jump Closure and Limit Uniformization\\in the Ideal Completion of the Turing Degrees}}
\author{Miara Sung}
\date{Source: arXiv:2608.25925v1 (CC BY 4.0)}
\begin{document}
\maketitle

\noindent\emph{Source and licence.} \textbf{Jump Closure and Limit Uniformization\\in the Ideal Completion of the Turing Degrees}. Version arXiv:2608.25925v1 (CC BY 4.0), distributed by arXiv under the Creative Commons Attribution 4.0 International licence, http://creativecommons.org/licenses/by/4.0/, as recorded in the arXiv licence metadata for this exact version and shown on the abstract page https://arxiv.org/abs/2608.25925v1. Full licence notice: \texttt{licenses/arxiv-2608.25925-CC-BY-4.0.txt}.\par\smallskip

\noindent\textit{Editorial scope.} Counted unit: the article's theorem thm:finite-block (source0.tex lines 685--698). The article's complete original proof of it is delivered with it (source0.tex lines 700--744, one \texttt{proof} environment of 45 lines). No mathematics is written, repaired or re-derived: the statement and the proof are the article's own lines, in the article's own notation, and no retained line is substituted or re-flowed.  Retained uncounted as the source-local context the proof needs: lines 310--327.  Nothing was omitted from any retained span, so the package carries no omission record.  No retained line contains a citation, so the wrapper carries no bibliography and no bibliography entry.  No retained line contains a figure environment, an image-inclusion command or drawing code, so no image is delivered and the package declares no asset dependency.  Editorial formatting only, in addition: an AMS \texttt{amsart} preamble that restates the article's own declarations and macros used by the retained lines, namely the environments \texttt{theorem} on separate counters, together with the standard packages those lines need.  The retained environments print in this document as Theorem 1 in order of appearance, whereas the article prints the same items with the numbering of its own full document.  The complete native source of the article as distributed in the arXiv e-print for this exact version is bundled unchanged as \texttt{sources/208605/source0.tex}.
\begin{theorem}[Least jump fixed point]
\label{thm:least-jump-fixed}
Starting from $\down\mathbf a$, the finite iterates of $\Gamma$ satisfy
\begin{equation}
\Gamma^n(\down\mathbf a)=\down\mathbf a^{(n)}
\label{eq:finite-iterates}
\end{equation}
for every $n<\omega$, and
\begin{equation}
\lfp_{\geq\down\mathbf a}(\Gamma)
=
\bigcup_{n<\omega}\down\mathbf a^{(n)}
=
\operatorname{Arith}(\mathbf a).
\label{eq:relative-lfp}
\end{equation}
The closure ordinal of this iteration is exactly $\omega$.
\end{theorem}

\begin{theorem}[Finite block theorem]
\label{thm:finite-block}
For every integer $m\geq1$, the closure ordinal of $\Theta_m$, starting from $\down\mathbf0$, is
\[
\omega m,
\]
and its least fixed point is
\begin{equation}
A_{\omega m}
=
\bigcup_{\beta<\omega m}\down\mathbf0^{(\beta)}.
\label{eq:A-omegam}
\end{equation}
\end{theorem}

\begin{proof}
The first block is the ordinary jump iteration:
\[
I_n=\down\mathbf j_n
\qquad(n<\omega),
\]
so $I_\omega=A_\omega$. If $m=1$, there is no gate at $\omega$, and \cref{thm:least-jump-fixed} gives the result.

Assume $m>1$. At $A_\omega$, jump closure adds nothing while the gate $U_\omega$ adds $\mathbf j_\omega$. Hence
\[
I_{\omega+1}=\down\mathbf j_\omega.
\]
Repeated jump closure gives
\[
I_{\omega+n+1}=\down\mathbf j_{\omega+n}
\qquad(n<\omega).
\]
At the next limit,
\[
I_{\omega\cdot2}=A_{\omega\cdot2}.
\]
If $2<m$, the $U_{\omega\cdot2}$ gate fires and yields
\[
I_{\omega\cdot2+1}=\down\mathbf j_{\omega\cdot2}.
\]
Continuing inductively, for every $1\leq k<m$,
\begin{equation}
I_{\omega k}=A_{\omega k}
\label{eq:block-boundary}
\end{equation}
and
\begin{equation}
I_{\omega k+n+1}=\down\mathbf j_{\omega k+n}
\qquad(n<\omega).
\label{eq:inside-block}
\end{equation}
At stage $\omega m$ we obtain $A_{\omega m}$.

This ideal is jump-closed. If $\mathbf x\leT\mathbf j_\beta$ for some $\beta<\omega m$, then
$\mathbf x'\leT\mathbf j_{\beta+1}$, and $\beta+1<\omega m$ because $\omega m$ is a limit ordinal. Thus $\Gamma(A_{\omega m})=A_{\omega m}$.

Every gate appearing in $\mathcal U_m$ corresponds to some $\omega k<\omega m$, and its limit degree $\mathbf j_{\omega k}$ already belongs to $A_{\omega m}$. Hence $\mathcal U_m(A_{\omega m})=A_{\omega m}$, so $A_{\omega m}$ is fixed.

No earlier stage is fixed. At a successor position strict jump progression forces movement. The only nonzero limit ordinals below $\omega m$ are the block boundaries $\omega k$ for $1\leq k<m$, and at each of those the corresponding uniformization gate forces movement. Hence the closure ordinal is exactly $\omega m$.
\end{proof}

\end{document}
